\documentclass{article}

%maths
\usepackage{mathtools}
\usepackage{amsmath}
\usepackage{amssymb}
\usepackage{amsfonts}

%tikzpicture
\usepackage{tikz}
\usepackage{scalerel}
\usepackage{pict2e}
\usepackage{tkz-euclide}
\usetikzlibrary{calc}
\usetikzlibrary{patterns,arrows.meta}
\usetikzlibrary{shadows}

%pgfplots
\usepackage{pgfplots}
\pgfplotsset{compat=newest}
\usepgfplotslibrary{statistics}
\usepgfplotslibrary{fillbetween}

%colours
\usepackage{xcolor}

\usepackage{nicefrac}

\begin{document}

\pgfplotsset{
	standard/.style={
		axis line style = thick,
		trig format=rad,
		enlargelimits,
		axis x line=middle,
		axis y line=middle,
		enlarge x limits=0.15,
		enlarge y limits=0.15,
		every axis x label/.style={at={(current axis.right of origin)},anchor=north west},
		every axis y label/.style={at={(current axis.above origin)},anchor=south east}
	}
}

\section{Relations And Functions}

\textbf{Relation} is a set of ordered pairs that can be represented by a diagram, graph, or sentence.\\
\\
The \textbf{Domain} is the set of all \textit{first} numbers of the ordered pairs in a relation.\\
The \textbf{Range} is the set of all \textit{second} numbers of the ordered pairs in a relation. \\
\\
A \textbf{Function} is when none of the domain values are repeated.\\
A \textbf{Vertical Line Test} is a test for functions.\\
\\
\subsection{Compositions of Functions}
A \textbf{composite function} is when a variable is replaced with another function.\\
\\
Example: f(g(x)) is read as: 'f of g of x'\\
\\
\textbf{Inverse Relations} is when you reverse ordered pairs of an initial relation.\\
The \textbf{identity function} is the line which if you folded the graph over on this line/crease it would line up both relations. Since for each replacement of x, the result is identical to x.

\begin{center}
	\begin{tikzpicture}
		\begin{axis}[standard, 
			xtick={-4,-3,-2,-1,0,1,2,3,4,5,6,7,8},
			ytick={-4,-3,-2,-1,0,1,2,3,4,5,6,7,8},
			samples=1000,
			xlabel={$x$},
			ylabel={$y$},
			xmin=-4,xmax=8,
			ymin=-4,ymax=8,
			clip=false]

			% horizontal line - a
			\addplot[
				<->,  
				thick, 
				domain=-4:8,
				samples=2
				] {x} node[anchor=south] {$y=x$};

			\fill (2,1) circle (2pt);
			\fill (1,2) circle (2pt);
			\fill (6,4) circle (2pt);
			\fill (4,6) circle (2pt);
			\fill (8,3) circle (2pt) node[anchor=south] {$R^{-1}$};
			\fill (3,8) circle (2pt) node[anchor=south] {$R$};






		\end{axis}
	\end{tikzpicture}
\end{center}

\end{document}
